\section{Conclusion}\label{sec:conclusion}

A broad disinformation index can be examined as a correlate of macroeconomic
risk without equating political-information scores with financial-message
contamination. The panel estimates a three-year cumulative-growth
slope of $\GrowthTwo$ percentage points at mean conditioning states and a
banking-stress interaction of $\FragilityTwo$. Quantile inference and
chronological prediction evaluate the scope of the evidence, including the
possibility that in-sample associations do not improve held-out forecasts.

The illustrative financing model supplies an increasing-differences result
under an explicitly assumed lending rule. Feedback can generate multiple
steady states for some parameters, but the empirical exercise does not
identify a trap threshold. Persistence comparative statics rank parameter
changes only under sign and effectiveness restrictions; they do not identify
policy cost-effectiveness. The remaining contribution is a reproducible
conditional analysis, with measurement, identification and forecasting
limitations made explicit. Financial-specific data and credible intervention
variation are the next steps needed to support causal policy conclusions.
